\documentclass[10pt]{article}
\usepackage{array, color, fancyhdr, multirow, longtable, enumitem}
\usepackage[top=0.5in, bottom=1in, left=0.75in, right=0.75in]{geometry}

\usepackage{hyperref}
\hypersetup{
    colorlinks=true,
    linkcolor=blue,
    filecolor=magenta,      
    urlcolor=cyan,
    bookmarks=true,
    bookmarksopen=true,
}
 
\urlstyle{same}

\usepackage{fontspec}
\setmainfont{Liberation Sans}[Scale=1]

\definecolor{lightgray}{gray}{0.8}
\newcolumntype{L}{>{\raggedleft\scriptsize}p{0.12\textwidth}}
\newcolumntype{R}{p{0.84\textwidth}}
\newcommand\VRule{\color{lightgray}\vrule width 0.5pt}

\setlength{\LTpre}{0pt}
%\setlength{\LTpost}{0pt}

\makeatletter
\renewcommand\section{\@startsection {section}{1}{\z@}%
                                   {-3.5ex \@plus -1ex \@minus -.2ex}%
                                   {1.0ex \@plus.2ex}%
                                   {\normalfont\large\bfseries}}% from \Large
\renewcommand\subsection{\@startsection{subsection}{2}{\z@}%
                                     {-3.25ex\@plus -1ex \@minus -.2ex}%
                                     {0.5ex \@plus .2ex}%
                                     {\normalfont\normalsize\bfseries}}% from \large
\makeatother

\setlength{\headheight}{0.5in}
\pagestyle{fancy}
\fancyhf{}

\rhead{Alan P. Boyle}
\lhead{\today}
\cfoot{\thepage}

\begin{document}
\thispagestyle{plain}

\begin{tabular*}{.95\textwidth}{l@{\extracolsep{\fill}}r}
& {\small 5B322 Medical Science I }\\
& {\small 1301 Catherine St. }\\
& {\small Ann Arbor, MI 48109-2218 }\\
\multirow{2}{*}{\textbf{\LARGE Alan P. Boyle}} & {\small apboyle@umich.edu }\\
& {\small \href{https://boylelab.org}{boylelab.org} }\\
\end{tabular*}
\hrule

\section*{Education}
\begin{tabular}{L!{\VRule}R}
2005--2009&{\bf Doctor of Philosophy,} Computational Biology and Bioinformatics\\
&{Duke University, Durham, NC}\\
\noalign{\medskip}
2001--2005&{\bf Bachelor of Science,} {\sl summa cum laude,} Biochemistry and Molecular Biology\\
&{\bf Bachelor of Science,} {\sl summa cum laude,} Computer Science\\
&{Mississippi State University, Starkville, MS}\\
\end{tabular}

\section*{Academic Appointments}
\begin{tabular}{L!{\VRule}R}
2025--present&{\bf Professor with tenure,} Department of Computational Medicine \& Bioinformatics\\
&{\bf Professor,} Department of Human Genetics\\
2023--present&{\bf Core Member,} Rogel Cancer Center\\
2020--2023&{\bf Affiliate Member,} Rogel Cancer Center\\
2021--present&{\bf Affiliate,} Michigan Neuroscience Institute\\
2016--present&{\bf Member,} Center for RNA Biomedicine\\
\noalign{\medskip}
2020--2025&{\bf Associate Professor with tenure,} Department of Computational Medicine \& Bioinformatics\\
&{\bf Associate Professor,} Department of Human Genetics\\
\noalign{\medskip}
2015--2020&{\bf Assistant Professor,} Department of Human Genetics\\
2014--2020&{\bf Assistant Professor,} Department of Computational Medicine \& Bioinformatics\\
&{University of Michigan, Ann Arbor, MI}\\
\noalign{\medskip}
2010--2014&{\bf Postdoctoral Scholar,} Genetics\\
&{Stanford University, Stanford, CA; Advisor: Dr. Michael Snyder}\\
\noalign{\medskip}
Spring 2010&{\bf Postdoctoral Associate,} Computational Biology\\
&{Duke University, Durham, NC; Advisor: Dr. Terrence S. Furey}\\
\end{tabular}

\section*{Scholarships, Fellowships, and Honors}
\begin{longtable}{L!{\VRule}R}
%2023&University of Michigan ‘Making a Difference’ Award from Office for Health Equity \& Inclusion\\
2022&Valuing our Own Award, Michigan Medicine\\
2019&Endowment for the Basic Sciences Teaching Award\\
2018&First Place in CAGI5 Regulation Saturation Challenge\\
2017&NSF CAREER Award\\
2016&Institutional nominee for W.M. Keck Foundation Medical Science Research Program\\
2016&Institutional nominee for Searle Scholar Award\\
2015--2017&Alfred P. Sloan Foundation Fellowship in Computational \& Evolutionary Molecular Biology\\
2013--2014&NIH Pathway to Independence Award (K99/R00) [1K99HG007356-01]\\
2012&AAAS/Science Program for Excellence in Science\\
2005--2008&NSF Graduate Research Fellowship\\
2005--2009&James B. Duke Fellowship\\
Summer 2004&{Mayo Clinic Summer Undergraduate Research Fellow}\\
2003&Barry M. Goldwater Memorial Scholarship\\
Summer 2003&{The Institute for Genomic Research (TIGR) Summer Fellow}\\
2001&Robert C. Byrd Honors Scholarship\\
2001&Mississippi State University Presidential Scholarship\\
2001&National Merit Scholarship
\end{longtable}

\section*{Grant Support}
\subsection*{Active}
\begin{longtable}{L!{\VRule}R}
2021--2026&U01 HG011952\hfill (PI: Boyle)\\
&NIH/NHGRI \hfill \\ %Total Costs: \textdollar3,829,187\\
&Predicting the impact of genomic variation on cellular states\\
&This project seeks to develop tools for interpretation of genomic variation on cellular state through modeling single cell data as part of the IGVF consortium.\\
\noalign{\medskip}
2022--2026&R01 GM144484 \hfill (PI: Boyle)\\
&NIH/NIGMS \hfill \\ %Total Costs: \textdollar1,560,000\\
&Mobile element derived chromatin looping variability in human populations\\
&This project seeks to study the impact of polymorphic LTR13 integrations on 3D chromatin conformation.\\
\noalign{\medskip}
2023--2028&UG3/UH3 NS132084 \hfill (Multi-PI: Mills, Boyle)\\
&NIH/OD \hfill \\ %Total Costs: \textdollar1,560,000\\
&Molecular and Computational Tools for Identifying Somatic Mosaicism in Human Tissues\\
&As part of the SMaHT consortium this project seeks to develop long-read methods to study somatic mosaicism in normal human tissues.\\
\noalign{\medskip}
2021--2026&K08 HL153799\hfill (PI: Denstaedt; Consultant)\\
&NIH/NHLBI \hfill \\ %Total Costs: \textdollar847,800\\
&Predisposition for Lung Injury in Sepsis Survival\\
&The goal of this project is to understand the biological mechanisms predisposing to these complications in order to prevent and treat them.\\
\noalign{\medskip}
2021--2026&R01 NS122165\hfill (PI: Castro; Co-I with Effort)\\
&NIH/NINDS \hfill \\ %Total Costs: \textdollar3,228,160\\
&Uncover the role of H3.3-G343R mutation in shaping the DNA damage response, anti-tumor immunity and mechanisms of resistance in glioma\\
&This project seeks to study pediatric high-grade gliomas with H3.3-G343R, ATRX, and TP53 inactivating mutations to understand the impact of H3.3-G343R on the tumor immune microenvironment.\\
\noalign{\medskip}
2022--2026&R01 CA260677\hfill (PI: Malek; Co-I with Effort)\\
&NIH/NCI \hfill \\ %Total Costs: \textdollar3,228,160\\
&The Biology of Mutant STAT6 in Follicular Lymphoma\\
&This project seeks to study STAT6 gene regulation in the context of B cell lymphoma.\\
\noalign{\medskip}
2023--2028&R01 NS099280\hfill (PI: Todd; Consultant)\\
&NIH/NINDS \hfill \\ %Total Costs: \textdollar3,228,160\\
&Hexanucleotide repeat translation in ALS and Frontotemporal Dementia\\
&This project seeks to study RAN translation in ALS and FTD at a hexonucleotide expansion in C9orf72.\\
\noalign{\medskip}
2024--2029&R01 DE032699 \hfill (PI: Brenner, Mills, Spector; Co-I with Effort)\\
&NIH/NIDCR \hfill \\ % Total Costs: \textdollar3,820,341\\
&Defining the Role of HPV Integration Structures in HNSCC Molecular Heterogeneity\\
&This proposal seeks to determine how the structure of HPV integration events influences chromatin accessibility and adjacent gene expression in HPV-positive oropharyngeal squamous cell carcinoma, with the goal of identifying pivotal driver integrations that can predict treatment outcomes and guide therapy decisions to improve patient survival and reduce morbidity.\\
\noalign{\medskip}
2026--2027&National Ataxia Foundation \hfill (PI: Burmeister, Boyle)\\
%&NIH/NIDCR \hfill \\ % Total Costs: \textdollar3,820,341\\
&FGF14 intronic GAA repeats in vestibular problems beyond SCA27B\\
&This project seeks to identify genetic properties of the SCA27B repeat that may represent subclinical expansions.\\
\end{longtable}

%\subsection*{Submitted}
%\begin{longtable}{L!{\VRule}R}
%2023--2028&R01 AG000000 \hfill (PI: McConnell; Co-I with Effort)\\
%&NIH/NIA \hfill \\ %Total Costs: \textdollar3,981,609\\
%&Somatic mutations and AD resilience\\
%\noalign{\medskip}
%2023--2028&R01 CA000000 \hfill (PI: Malek; Co-I with Effort)\\
%&NIH/NCI \hfill \\ %Total Costs: \textdollar1,950,000\\
%&The Biology of Mutant STAT6 in Follicular Lumphoma\\
%\noalign{\medskip}
%2022--2024&R21 CA261641 \hfill (PI: Boyle)\\
%&NIH/NCI \hfill \\ %Total Costs: \textdollar599,846\\
%&Direct capture of complete HPV integration sites using long-read sequencing\\
%\noalign{\medskip}
%2022--2024&R21 HG012849 \hfill (PI: Zhou; Co-I with Effort)\\
%&NIH/NHGRI \hfill \\ %Total Costs: \textdollar416,696\\
%&A multi-omics framework for detection and functional analysis of germline and somatic transposable elements in human tissues\\
%\end{longtable}

\subsection*{Completed}
\begin{longtable}{L!{\VRule}R}
2013--2017&R00 HG007356 Pathway to Independence Award (K99/R00) \hfill (PI: Boyle)\\
&NIH/NHGRI \hfill \\ %Total Costs: \textdollar987,771\\
&Global Discovery and Validation of Functional Regulatory Elements\\
&This project seeks to extend current assays demonstrating function of genomic regions into an equivalent genome-wide assay.\\
\noalign{\medskip}
2015--2017&FG-2015-65465 \hfill (PI: Boyle)\\
&Alfred P. Sloan Foundation \hfill \\ %Total Costs: \textdollar50,000\\
&Fellowship in Computational \& Evolutionary Molecular Biology\\
\noalign{\medskip}
2016--2020&R01 HL130705 \hfill (PI: Willer; Co-I with Effort)\\
&NIH/NHLBI \hfill \\ %Total Costs: \textdollar2,784,005\\
&Large-scale human genetics to understand molecular mechanisms of atrial fibrillation and related traits\\
&This project seeks to provide new insights into atrial fibrillation mechanisms through whole-genome screening.\\
\noalign{\medskip}
2017--2018&Eleanor and Larry Jackier U-M/Technion and Weizmann Collaborative Research Grant \\
& \hfill (co-PI: Boyle, Mandel-Gutfreund)\\
&Michigan - Israel Partnership for Research \& Education \hfill \\ %Total Costs: \textdollar50,000\\
&Identifying novel disease related mutations in the genomic environments around Trascription Factor binding sites \\
&The goal of this project is to identify variants in the proximity of TF binding sites that have an indirect effect on their binding.\\
\noalign{\medskip}
2017--2024&R35 HL135824 \hfill (PI: Willer; Co-I with Effort)\\
&NIH/NHLBI \hfill \\ %Total Costs: \textdollar4,650,000\\
&Using Genetics to Inform Mechanism of Cardiovascular Disease\\
&The goal of this project is to uncover novel genetic discoveries and biological mechanisms underlying association with devastating cardiovascular diseases.\\
\noalign{\medskip}
2019&NVIDIA GPU Grant  \hfill (PI: Boyle)\\
&NVIDIA Corporation\hfill \\ %In-Kind Value: \textdollar3,600\\
\noalign{\medskip}
2017--2022&DBI-1651614 \hfill (PI: Boyle)\\
&NSF/BIO/DBI \hfill \\ %Total Costs: \textdollar979,984\\
&CAREER: Conservation of cohesin-containing cis regulatory modules in the human and mouse lineages\\
&The goal of this project is the study of the turnover of cohesin binding sites in the human and mouse genomes.\\
\noalign{\medskip}
2022--2022&R21 HG011493 S1\hfill (Multi-PI: Boyle, Mills)\\
&NIH/NIA \hfill \\ %Total Costs: \textdollar412,928\\
&New technologies for accurate capture and sequencing of repeat-associated regions (Supplement)\\
&This project seeks to map mobile elements in a set of Alzheimer's samples.\\
\noalign{\medskip}
2019--2022&Precision Health Investigators Award\hfill (co-PI: Todd, Boyle, Mills)\\
&University of Michigan \hfill \\ %Total Costs: \textdollar300,000\\
&Short Tandem repeats in precision health and human disease\\
&The goal of this project is to develop any assay to measure STRs in human genomes and develop bioinformatic tools to predict STR expansions from genotypes.\\
\noalign{\medskip}
2022&NVIDIA GPU Grant  \hfill (PI: Boyle)\\
&NVIDIA Corporation\hfill \\ %In-Kind Value: \textdollar30,000\\
\noalign{\medskip}
2021--2022&Cancer Center Discovery\hfill (PI: Boyle)\\
&University of Michigan \hfill \\ %Total Costs: \textdollar75,000\\
&Direct capture of complete HPV integration sites using long-read sequencing\\
&This project seeks to develop methods to capture of complete HPV integration events in the human genome.\\
\noalign{\medskip}
2020--2023&W81XWH2010336 \hfill (PI: Aguilar; Co-I with Effort)\\
&DoD/Army \hfill \\ %Total Costs: \textdollar1,081,608\\
&Understanding \& Enhancing the Regenerative Capacity of Skeletal Muscle to Trauma by Targeting Muscle-Nerve Synergy\\
&This project seeks to study the single cell chromatin and RNA landscape in skeletal muscle repair.\\
\noalign{\medskip}
2020--2023&R21 HG011493\hfill (Multi-PI: Boyle, Mills)\\
&NIH/NHGRI \hfill \\ %Total Costs: \textdollar412,928\\
&New technologies for accurate capture and sequencing of repeat-associated regions\\
&This project seeks to map mobile elements in a trio of cell lines and develop technologies for improving this mapping.\\
\noalign{\medskip}
2018--2024&R01 HD093570 \hfill (PI: Bielas; Co-I with Effort)\\
&NIH/NICHD \hfill \\ %Total Costs: \textdollar2,304,265\\
&Genetic Diagnosis of Neurodevelopmental Disorders in India\\
&This study will establish whole-exome sequencing to study mendelian genetic disorders at the All India Institute of Medical Sciences.\\
\noalign{\medskip}
2022--2024&Michigan Alzheimer's Disease Center Developmental Project\hfill (PI: Zhou; Consultant)\\
&University of Michigan \hfill \\ %Total Costs: \textdollar50,000\\
&Explore the functional impact of transposable elements in Alzheimer’s disease and related dementias\\
&This project seeks to explore the connection between the somatic transposable elements in the human genome and Alzheimer’s disease and related dementias. \\
\noalign{\medskip}
2022--2025&R21 CA2578964\hfill (PI: Boyle)\\
&NIH/NCI \hfill \\ %Total Costs: \textdollar561,860\\
&High-throughput inverted reporter assay for characterization of silencers and enhancer blockers\\
&This project seeks to develop tools for the study of negative regulatory elements in cancer development.\\
\noalign{\medskip}
2021--2025&R01 HD104680\hfill (PI: Hammoud; Co-I with Effort)\\
&NIH/NICHD \hfill \\ %Total Costs: \textdollar1,663,876\\
&Sperm Chromatin: Implications on organismal development and fertility\\
&This project seeks to explore protamine chromatin structure in mouse sperm.\\
\noalign{\medskip}
2017--2025&U24 HG009293 \hfill (Multi-PI: Boyle, Cherry)\\
&NIH/NHGRI \hfill \\ %Total Costs: \textdollar2,171,753\\
&RegulomeDB: A Resource for the Human Regulome\\
&This project seeks to expand and support a RegulomeDB, a database for prioritizing and predicting functional variants in the human genome.\\
\noalign{\medskip}
2022--2025&Taubman Institute Innovation Projects\hfill (co-PI: Todd, Boyle, Mills)\\
&University of Michigan \hfill \\ %Total Costs: \textdollar300,000\\
&Short Tandem repeats in precision health and human disease\\
&The goal of this project is to develop any assay to measure STRs in human genomes and develop bioinformatic tools to predict STR expansions from genotypes.\\
\end{longtable}


\section*{Professional Service}
\subsection*{Service}
\begin{longtable}{L!{\VRule}R}
2022--current&University of Michigan Biomedical Research Council (BMRC) [Standing Member; co-Chair 2026--current]\\
2026--current&Michigan Neuroscience Institute (MNI) Executive Committee\\
2014--2018, 2025--current&DCM\&B Admissions Committee [Chair 2025--current]\\
2024--current&DHG Faculty Development Committee\\
2023--current&R01 Bootcamp Medical School Cohort Coach\\
2023--current&Somatic Mosaicism across Human Tissues (SMaHT) Consortium Steering Committee\\
2021--current&Impact of Genomic Variation on Function (IGVF) Consortium Steering Committee\\
2018--current&Lab Safety Liaison for DCM\&B\\
2015, 2016, 2025, 2026&DCM\&B Retreat Planing Committee Chair (including 1st annual)\\
2020--2022&DHG M.S. Admissions Committee\\
%2018--current&DCM\&B Diversity, Equity, \& Inclusion Committee [Ally/Chair 2018--2020]\\
2017--2024&DCM\&B Preliminary Exam Abstract Review Committee (PARC) [Chair 2018--2022]\\
2019--2020&DHG Ph.D. Admissions Committee\\
2017--2020&DHG Faculty Recruitment and Promotions Committee\\
2016--2020&DCM\&B Seminar Series Committee [Chair]\\
2018--2019&Cellular and Molecular Biology Admissions Committee\\
2017--2019&EBS Faculty IT Committee\\
2016--2019&DCM\&B Faculty Recruitment Committee\\
2015--2017&DHG Computational Support Committee\\ 
2008--2009&Duke Computational Biology \& Bioinformatics student committee\\ 
\end{longtable}

\subsection*{Memberships}
\begin{longtable}{L!{\VRule}R}
2018--current&Member, American Society of Human Genetics (ASHG)\\
2013--current&Member, International Society for Computational Biology (ISCB)\\
2012--current&Member, American Association for the Advancement of Science (AAAS)\\
2005--current&Member, Gamma Sigma Delta Agricultural Honor Society\\
\end{longtable}

\subsection*{Manuscript Reviewing Activity}
\begin{longtable}{L!{\VRule}R}
Since 2009&{\sl Ad hoc} reviewer \href{https://publons.com/researcher/1152609/alan-boyle/}{(>100 verified reviews)} for the journals: {\sl Science}, {\sl Nature Biotechnology}, {\sl Nature Genetics}, {\sl Genome Research}, {\sl Genome Biology}, {\sl Nature Neuroscience}, {\sl Nature Communications}, {\sl Nature Protocols}, {\sl Bioinformatics}, {\sl Nucleic Acids Research}, {\sl BMC Biology}, {\sl BMC Bioinformatics}, {\sl PLOS Computational Biology}, {\sl Oncotarget}, {\sl Scientific Reports}, {\sl Atherosclerosis}, {\sl BioEssays}, {\sl Gene}\\
2026&Academic Editor, PLOS One\\
2026&Program Committee, Proceedings, ISMB/ECCB\\
2023&Program Committee, Genome Sequence Analysis, ISMB/ECCB\\
2023&Program Committee, Biomedical Informatics, ISMB/ECCB\\
2018, 2020&Program Committee, Comparative and Functional Genomics, ISMB/ECCB\\
2018, 2019&Program Committee, Studies of Phenotypes and Clinical Applications, ISMB/ECCB\\
2019&Program Committee, General Computational Biology, ISMB/ECCB\\
2017&Program Committee, Regulatory Genomics Special Interest Group Meeting (RegGenSIG), ISMB/ECCB\\
2015--2018&Program Committee, Great Lakes Bioinformatics and Canadian Computational Biology Conference (GLBIO/CCBC)\\
2015--2016&Program Committee, Algorithms for Computational Biology (ALCOB)\\
2013--2016&Program Committee, Gene Regulation and Transcriptomics, ISMB/ECCB\\
2012--2015&DNA Day Essay Contest Detailed Review Judge for ASHG\\
2012&Distinguished contributor as a leading reviewer for the journal {\sl Bioinformatics}\\
\end{longtable}

\subsection*{Grant Reviewing Activity}
\begin{longtable}{L!{\VRule}R}
2026&NIH Study Section BDMA - Biodata Management and Analysis Study Section (Ad Hoc)\\
2026&NIH Study Section ZRG1 IIDB Z (70) - Stage I New Innovator Award Program (Mail-in)\\
2026&NIH Study Section ZRG1 MGG-F (90) S - Special Topics in Genomics, Genetic Evolution, and Technology Development (Ad Hoc)\\
2025&NIH Study Section BDMA - Biodata Management and Analysis Study Section (Ad Hoc)\\
2024&NIH Study Section ZRG1 BBBT-D (50) - PAR Panel: Biodata Management and Common Fund Data Sets (Ad Hoc and co-Chair)\\
2024&NIH Study Section BDMA - Biodata Management and Analysis Study Section (Ad Hoc)\\
2023&NSF Review Panel - Molecular and Cellular Biosciences (MCB) - Genetic Mechanisms (Ad Hoc)\\
2023&NIH Study Section - Multi-Omics of Health and Disease – Data Analysis and Coordination Center\\
2023&NIH Study Section GVE - Genetic Variation and Evolution Study Section (Ad Hoc)\\
2022&NIH Study Section ZRG1 ISB-S (57) - Academic-Industrial Partnerships for Translation of Technologies for Diagnosis and Treatment\\
2022&NASA Study Section E.11 Space Biology: Animal Studies - Omics Systems [21SBAS-OmisSys] (Ad Hoc)\\
2020&NIH/NIMH Study Section ZMH1 ERB-C (08) - Fine-Mapping Genome-Wide Associated Loci to Identify Proximate Causal Mechanisms of Serious Mental Illness\\
2019&NIH/NIMH Study Section ZMH1 ERB-C (01) - PsychENCODE: Non-Coding Functional Elements in the Human Brain and Their Role in the Development of Psychiatric Disorders\\
2018--2019&University of Michigan internal review for Searle Scholars Program\\
2015&UK Medical Research Council (RCUK MRC) - Methodology Research Panel (Ad Hoc)\\
2015&UK Biotechnology and Biological Sciences Research Council (RCUK BBSRC) (Ad Hoc)\\
2015&Michigan Institute for Clinical \& Health Research (MICHR) Postdoctoral Translational Scholars Program (Ad Hoc)\\
\end{longtable}

\section*{Teaching and Mentorship}
\section*{Training Programs}
\begin{longtable}{L!{\VRule}R}
2024--present&{\bf Member,} Systems \& Integrative Biology Training Grant (SIB)\\
2021--present&{\bf Member,} Biomedical Informatics and Data Science Training Program (BIDS-TP)\\
2017--present&{\bf Member,} Cellular and Molecular Biology Program\\
2015--present&{\bf Member,} Genome Science Training Program (GSTP)\\
2015--present&{\bf Member,} Michigan Predoctoral Training Program in Genetics (GTP)\\
2014--present&{\bf Member,} Program in Biomedical Sciences\\
2014--present&{\bf Member,} Bioinformatics Training Program\\
\end{longtable}

\subsection*{Teaching \textmd{\scriptsize(F = Fall Term, W = Winter Term, S = Summer Term)}}
\begin{longtable}{L!{\VRule}R}
W19, W20, W21, W22, W23, W24, W25, W26&Bioinformatics Concepts and Algorithms (BIOINF 529) [Course Director]\\
F15, F16, F17, F18, F19, F20, F21&Gene Structure and Regulation (HUMGEN 541) [3 lectures + 2 discussions / yr.]\\
F19, F22&Research Responsibility and Ethics (PIBS 503) [1 discussion / yr.]\\
F21, W22, F24, W25&Genetics Student Seminar (HUMGEN 821/822) [Mentor]\\
F17, F18&Experimental Genetics Systems (HUMGEN 632) [Course Director]\\
F15, W16, F16, W17, F17, W18, F18&Bioinformatics Journal Club (BIOINF 602/603) [Course Director F18]\\ 
S17, S18&Introduction to Biocomputing Bootcamp (BIOSTAT/BIOINF/HUMGEN 606) [2 full days / yr.]\\
F15, F16, F17&Introduction to Bioinformatics \& Computational Biology (BIOINF 527) [2 lectures + 3 labs / yr.]\\
S15, S16, S17&Basic Biology for Graduate Students with Quantitative Training (BIOINF 523) [2 lectures / yr.]\\
F03&{Lab TA for Isotopes Tech I (MS. State, BCH 4414)}\\
\end{longtable}

\subsection*{Guest Lectures / Panels}
\begin{longtable}{L!{\VRule}R}
2018--2019&Lecturer, REU Site: Mathematical and Theoretical Biology Institute (MTBI), Arizona State University (NSF1757968) [2 days]\\
2017&Panel member, U. Michigan ``New Faculty Orientation to Corporate \& Foundation Relations" [~70 attendees]\\
2016&Experimental Genetics Systems (HUMGEN 632) [1 discussion]\\
2014&Panel member, BIOINF 527 ``Challenges in Biology, Biomedicine, Data \& Analysis"\\
2010&Co-taught Cold Spring Harbor Systems Biology Pre-meeting Workshop\\
2009&Duke student panelist for ``How to prepare for and get into graduate school"\\
2008&Taught Duke mini-course on Genome Browsers \& Databases\\
\end{longtable}

\subsection*{Mentorship} 
\subsubsection*{Ph.D. Students (n=21)} 
\begin{longtable}{L!{\VRule}R} 
2026--&Kobe Howcroft (Ph.D. Student, Bioinformatics, University of Michigan)\\
\noalign{\medskip}
2024--current&Ingrid Flaspohler (Ph.D. Student, Bioinformatics, University of Michigan)\\
	&{\sl NIH Biomedical Informatics and Data Science Training Program (T32)}\\ 
	&{\sl Rackham Graduate Student Research Grant (pre-candidate)}\\ 
	&{\sl Rackham Graduate Student Research Grant (candidate)}\\ 
\noalign{\medskip}
2024--current&Steve Losh (Ph.D. Student, Bioinformatics, University of Michigan)\\
\noalign{\medskip}
2024--current&Sowmya Srinivasan (Ph.D. Student, Genetics and Genomics, University of Michigan)\\
\noalign{\medskip}
2022--current&Katarina Pavlovic (Ph.D. Student, Bioinformatics, University of Michigan)\\
	&{\sl Rackham Graduate Student Research Grant (pre-candidate)}\\ 
	&{\sl Rackham Graduate Student Research Grant (candidate)}\\ 
\noalign{\medskip}
2022--current&Rintsen Sherpa (Ph.D. Student, Bioinformatics, University of Michigan)\\
	&{\sl Rackham Graduate Student Research Grant (candidate)}\\ 
\noalign{\medskip}
2021--2025&Kinsey Van Deynze (Ph.D. Student, Bioinformatics, University of Michigan)\\
	&{\sl NIH Genome Science Training Program (T32)}\\ 
	&{\sl Rackham Graduate Student Research Grant (pre-candidate)}\\ 
	&{\sl Rackham Graduate Student Research Grant (candidate)}\\ 
\noalign{\medskip}
2020--2025&Andrea Valenzuela (Ph.D. Student, Chemical Biology, University of Michigan)\\
	&{\sl NIH Cellular Biotechnology Training Program (T32)}\\ 
\noalign{\medskip}
2020--2026&Breanna McBean (Ph.D. Student, Genetics and Genomics, University of Michigan)\\
	&{\sl Joint M.S. in Bioinformatics, University of Michigan}\\
	&{\sl NIH Genome Science Training Program (T32)}\\ 
	&{\sl Rackham Graduate Student Research Grant (pre-candidate)}\\ 
	&{\sl Rackham Graduate Student Research Grant (candidate)}\\ 
\noalign{\medskip}
2020--2025&Camille Mumm (Ph.D. Student, Genetics and Genomics, University of Michigan)\\
	&{\sl Joint M.S. in Bioinformatics, University of Michigan}\\
	&{\sl NIH Genome Science Training Program (T32)}\\ 
	&{\sl Rackham Graduate Student Research Grant (pre-candidate)}\\ 
	&{\sl Rackham Graduate Student Research Grant (candidate)}\\ 
	&{\sl Rackham Pre-doctoral Fellowship}\\ 
\noalign{\medskip}
2018--2024&Bradley Crone (Ph.D. Student, Bioinformatics, University of Michigan)\\
	&{\sl Rackham Graduate Student Research Grant (candidate)}\\ 
\noalign{\medskip}
2017--2023&Melissa Englund (Ph.D. Student, Genetics and Genomics, University of Michigan)\\ 
	&{\sl NIH Human Genetics Training Program (T32)}\\ 
	&{\sl Rackham Graduate Student Research Grant (candidate)}\\ 
\noalign{\medskip}
2018--2023&Nanxiang (Samuel) Zhao (Ph.D. Student, Bioinformatics, University of Michigan)\\ 
	&{\sl Rackham Graduate Student Research Grant (pre-candidate)}\\ 
	&{\sl Rackham Graduate Student Research Grant (candidate)}\\ 
\noalign{\medskip}
2016--2018&Haley Amemiya (Ph.D. Student, Cellular and Molecular Biology, University of Michigan)\\ 
	&{\sl Joint M.S. in Bioinformatics, University of Michigan}\\
	&{\sl NIH Cellular \& Molecular Biology Training Program (T32)}\\ 
	&{\sl NIH Cellular Biotechnology Training Program (T32) (Declined)}\\ 
	&{\sl PIBS Excellence in Service Award}\\ 
	&{\sl Rackham Graduate Student Research Grant (pre-candidate)}\\ 
	&{\sl Rackham Graduate Student Research Grant (candidate)}\\ 
	&{\sl Maas Professional Development Award}\\ 
	&{\sl Rackham Graduate School Scholar-Activist Award}\\
\noalign{\medskip}
2016--2020&Shriya Sethuraman (Ph.D. Student, Bioinformatics, University of Michigan)\\ 
\noalign{\medskip}
2016--2023&Christopher Castro (Ph.D. Student, Bioinformatics, University of Michigan)\\ 
	&{\sl NIH Bioinformatics Training Program (T32)}\\ 
	&{\sl Rackham Merit Fellow}\\ 
	&{\sl Rackham Graduate Student Research Grant (pre-candidate)}\\
	&{\sl Rackham Graduate Student Research Grant (candidate)}\\ 
	&{\sl Global Research Engagement Opportunity Fellowship}\\ 
\noalign{\medskip}
2017--2022&Ningxin Ouyang (Ph.D. Student, Bioinformatics, University of Michigan)\\ 
	&{\sl Rackham Graduate Student Research Grant (candidate)}\\ 
\noalign{\medskip}
2016--2021&Shengcheng Dong (Ph.D. Student, Bioinformatics, University of Michigan)\\ 
	&{\sl Rackham Graduate Student Research Grant (candidate)}\\ 
\noalign{\medskip}
2015--2021&Torrin McDonald (Ph.D. Student, Genetics and Genomics, University of Michigan)\\ 
	&{\sl NIH Human Genetics Training Program (T32)}\\ 
	&{\sl Rackham Graduate Student Research Grant (pre-candidate)}\\ 
	&{\sl Rackham Graduate Student Research Grant (candidate)}\\ 
\noalign{\medskip}
2015--2017&Greg Farnum (Ph.D. Student, Cellular and Molecular Biology, University of Michigan)\\ 
\noalign{\medskip}
2015--2020&Sierra Nishizaki (Ph.D. Student, Genetics and Genomics, University of Michigan)\\ 
	&{\sl Joint M.S. in Bioinformatics, University of Michigan}\\
	&{\sl NIH Genome Science Training Program (T32)}\\ 
	&{\sl Rackham Merit Fellow}\\ 
	&{\sl Rackham Summer Award}\\ 
	&{\sl Rackham Graduate Student Research Grant (candidate)}\\ 
\end{longtable}

\subsubsection*{M.S. Students (n=7)} 
\begin{longtable}{L!{\VRule}R} 
2025--2025&Christine Geng (M.S. Student, Bioinformatics, University of Michigan)\\
2023--2023&Hawra Aljawad (M.S. Student, Chemical Engineering, University of Michigan)\\
2023--2023&Xinyi Liu (M.S. Student, Bioinformatics, University of Michigan)\\
2022--2023&Emily Pogson (M.S. Student, Genetics and Genomics, University of Michigan)\\
2019--2020&Monica Holmes (M.S. Student, Bioinformatics, University of Michigan)\\
2017--2018&Nanxiang (Samuel) Zhao (M.S. Student, Bioinformatics, University of Michigan)\\ 
2015--2017&Ningxin Ouyang (M.S. Student, Bioinformatics, University of Michigan)\\ 
\end{longtable}

\subsubsection*{Additional Graduate Rotation Students (n=17)}
\begin{longtable}{L!{\VRule}R}
2024&Karan Smith (Rotation Student, Cellular and Molecular Biology, University of Michigan)\\
2024&Jun Sik Yun (Rotation Student, Genetics and Genomics, University of Michigan)\\
2024&Jeremy Chen (Rotation Student, Bioinformatics, University of Michigan)\\
2023&Rosina Carr (Rotation Student, Bioinformatics, University of Michigan)\\
2023&Connor Ward (Rotation Student, Medical Science Training Progran, University of Michigan)\\
2022&Brandt Bessell (Rotation Student, Bioinformatics, University of Michigan)\\
2022&Xiaomeng Du (Rotation Student, Bioinformatics, University of Michigan)\\
2022&Mahnoor Gondal (Rotation Student, Bioinformatics, University of Michigan)\\
2022&Xin Li (Rotation Student, Biological Chemistry, University of Michigan)\\
2022&Bohan Chen (Rotation Student, Cell and Developmental Biology, University of Michigan)\\
2021&Amelia Lauth (Rotation Student, Cellular and Molecular Biology, University of Michigan)\\
2019&Margarita Brovkina (Rotation Student, Cellular and Molecular Biology, University of Michigan)\\
2018&Steve Ho (Rotation Student, Human Genetics, University of Michigan)\\
2018&Matthew Pun (Rotation Student, Medical Science Training Progran, University of Michigan)\\
2017&Amanda Moccia (Rotation Student, Human Genetics, University of Michigan)\\
2017&Stephen Carney (Rotation Student, Human Genetics, University of Michigan)\\
2016&Tingyang Li (Rotation Student, Bioinformatics, University of Michigan)\\
\end{longtable}

\subsubsection*{Postdoctoral Fellows (n=4)}
\begin{longtable}{L!{\VRule}R}
2026--current&Kinsey Van Deynze (University of Michigan)\\
2023--2025&Melissa Englund (University of Michigan)\\
2022--2025&Torrin McDonald (University of Michigan)\\
2021--2022&Shengcheng Dong (University of Michigan)\\
\end{longtable}

\subsubsection*{Non-student Lab Volunteers (n=2)}
\begin{longtable}{L!{\VRule}R}
2019--2021&Greg Farnum (University of Michigan)\\
2018--2019&Monica Holmes (Postbac, University of Michigan)\\
\end{longtable}


\subsubsection*{Undergraduate and High School Students (n=27)}
\begin{longtable}{L!{\VRule}R}
2026--current&Pria Chandarana (Undergraduate, Neuroscience, University of Michigan)\\
2026--current&Cecelia Crandall (Undergraduate, Cognitive Science, University of Michigan)\\
2026--current&Boutheine Teyeb (Undergraduate, Computer Science, University of Michigan)\\
2026--2026&Aanshi Agrawal (Undergraduate, Molecular, Cellular, and Developmental Biology, University of Michigan)\\
2025--2026&Krrish Thakker (Undergraduate, Biochemistry, University of Michigan)\\
2023--2024&Kateri Darr (Undergraduate, Computer Science, University of Michigan)\\
2023--2023&Mason Miller (Undergraduate, Computer Science, University of Michigan)\\
2022--2024&Summer Ann (Undergraduate, Neuroscience, University of Michigan)\\
2022--2025&Kobe Howcroft (Undergraduate, Computer Science, University of Michigan)\\
2021--2024&Preston Parana (Undergraduate, UROP Molecular, Cellular, and Developmental Biology, University of Michigan)\\
	&{\sl UROP Blue Ribbon Award}\\ 
2021--2022&Julia Tweadey (Undergraduate, LSA Honors Program, Life Science Informatics, University of Michigan)\\
2021&Aryn Booker (Undergraduate, UROP Molecular, Cellular, and Developmental Biology, University of Michigan)\\
	&{\sl UROP Blue Ribbon Award}\\ 
2020&Marcela Alcaide Aligio (Undergraduate, SROP, Hunter College CUNY)\\
2019--2020&David Wang (Undergraduate, UROP Computer Science, University of Michigan)\\
2019--2020&Jack Lu (Undergraduate, UROP Computer Science, University of Michigan)\\
2019--2020&Diana Davis (Undergraduate, Neuroscience and German, University of Michigan)\\
2019&Sheila Rasouli (Undergraduate, Neuroscience, University of Toronto)\\
2019&Vibhasri Davuluri (High School, Girls Who Code Summer Intern)\\
2016--2019&Cody Morterud (Undergraduate, UROP Computer Science / Honors Capstone, University of Michigan)\\
2016--2017&Colten Williams (Undergraduate, UROP Computer Science, University of Michigan)\\
2016--2017&Courtney Asman (Undergraduate, Neuroscience, University of Michigan)\\
2014--2017&Maxwell Spadafore (Undergraduate, LS\&A Honors Informatics, University of Michigan)\\
2013--2014&Natalie Ng (High School, Stanford Institutes of Medicine Summer Research)\\
2013--2014&Dana Wyman (Undergraduate, Biology, Stanford University)\\
2013&Justin Young (High School, Stanford Institutes of Medicine Summer Research)\\
2012&Melanie Connick (Undergraduate, Biology, University of New Mexico)\\
2012&Edward Dai (Undergraduate, Computer Science, Stanford University)\\
\end{longtable}

\subsubsection*{Doctoral Thesis Committees (n=53)}
\begin{longtable}{L!{\VRule}R}
2026--current&Connor Davison (Genetics and Genomics, University of Michigan, Committee Member)\\
2026--current&Kobe Howcroft (Bioinformatics, University of Michigan, Chair)\\
2025--current&Brandt Bessell (Bioinformatics, University of Michigan, Committee Member)\\
2024--current&Weizhou Qian (Bioinformatics, University of Michigan, Committee Member)\\
2024--current&Jinhao Wang (Bioinformatics, University of Michigan, Committee Member)\\
2024--current&Brandon Klein (Medicinal Chemistry, University of Michigan, Committee Member)\\
2024--current&Bohan Chen (Bioinformatics, University of Michigan, Committee Member)\\
2024--current&Sowmya Srinivasan (Genetics and Genomics, University of Michigan, co-Chair)\\
2024--current&Ingrid Flaspohler (Bioinformatics, University of Michigan, Chair)\\
2024--current&Steve Losh (Bioinformatics, University of Michigan, Chair)\\
2024--current&Lu Lu (Bioinformatics, University of Michigan, Committee Member)\\
2023--current&Linghua Jiang (Bioinformatics, University of Michigan, Committee Member)\\
2023--current&Elysia Chou (Bioinformatics, University of Michigan, Committee Member)\\
2023--current&Rebecca McAvoy (Molecular, Cellular, and Developmental Biology, University of Michigan, Committee Member)\\
2023--current&Chinmay Raut (Bioinformatics, University of Michigan, Committee Member)\\
2022--current&Katarina Pavlovic (Bioinformatics, University of Michigan, Chair)\\
2022--current&Rintsen Sherpa (Bioinformatics, University of Michigan, Chair)\\
2022--current&Emily Peirent (Neuroscience, University of Michigan, Committee Member)\\
2025--current&Monica Holmes (Bioinformatics, University of Michigan, Committee Member)\\
&{\sl Single-Molecule Long-read Sequencing for Resolution of Structurally Complex Genomic Loci and Accurate Phenotype Prediction}\\
2024--2026&Matthew Hodgman (Bioinformatics, University of Michigan, Committee Member)\\
&{\sl Raising the Standard of Machine Learning Transparency and Utility in Cardiovascular Medicine}\\
2020--2026&Breanna McBean (Genetics and Genomics, University of Michigan, co-Chair)\\
&{\sl Targeting Resistance and Metastasis: Advancing Local and Systemic Therapeutic Strategies for Aggressive Breast Cancers}\\
2021--2026&Mashiat Rabbani (Genetics and Genomics, University of Michigan, Committee Member)\\
&{\sl A Genome-scale Program for Building the Sperm Epigenome}\\
2024--2025&Lingrui Cai (Bioinformatics, University of Michigan, Committee Member)\\
&{\sl Advancing Medical Image and Report Analysis with Artificial Intelligence: Applications in Inflammatory Bowel Disease and Traumatic Brain Injury}\\
2021--2025&Kinsey Van Deynze (Bioinformatics, University of Michigan, Chair)\\
&{\sl Sequence Analysis of Tandem Repeats Utilizing 3rd Generation Sequencing Technologies}\\
2020--2025&Andrea Valenzuela (Chemical Biology, University of Michigan, co-Chair)\\
&{\sl AI-Driven Workflows Integrating Strain Diversity and Host Metabolic-Genetic Variation for Tuberculosis Treatment Prioritization}\\
2021--2025&Zhenhao Zhang (Bioinformatics, University of Michigan, Committee Member)\\
&{\sl Decoding Genome Regulation through Deep Learning Models}\\
2020--2025&Steve Ho (Genetics and Genomics, University of Michigan, Committee Member)\\
&{\sl Applied Machine Learning for Big Data Genomics}\\
2020--2025&Camille Mumm (Genetics and Genomics, University of Michigan, Chair)\\
&{\sl Exploring Repetitive Elements in Neurodegenerative Disease Using Targeted Long-Read Sequencing}\\
2022--2025&Kaiwen Deng (Bioinformatics, University of Michigan, Committee Member)\\
&{\sl Advancing Computational Models and Algorithms for the Analysis of Single-Cell and Neural Data}\\
2022--2024&Franco Tavella (Biophysics, University of Michigan, Committee Member)\\
&{\sl Robustness and Tunability of Biological Oscillations}\\
2018--2024&Bradley Crone (Bioinformatics, University of Michigan, Chair)\\
&{\sl Computational Methods in Functional Prioritization of Polygenic Risk Score Models}\\
2021--2024&Wenjin Gu (Bioinformatics, University of Michigan, Committee Member)\\
&{\sl Development of Viral Integration Analysis Technologies for Virus-Associated Cancer Research}\\
2018--2023&Rucheng Diao (Bioinformatics, University of Michigan, Committee Member)\\
&{\sl Local Chromatin Environments Shape Transcription and Adaptive Immunity in Bacteria}\\
2021--2023&Zijun Gao (Bioinformatics, University of Michigan, Committee Member)\\
&{\sl Advance Machine Learning and Image Analysis Methods for Clinical Decision Support in Cardiovascular and Pulmonary Diseases}\\
2018--2023&Nanxiang (Samuel) Zhao (Bioinformatics, University of Michigan, Chair)\\ 
&{\sl Decoding Regulatory Variants with Computational Methods in Non-coding Regions of the Human Genome}\\
2020--2023&Ashley Melnick (Cellular and Molecular Biology, University of Michigan, Committee Member)\\
&{\sl Cdc73 Protects Notch-Induced Leukemia Cells From DNA Damage and Mitochondrial Stress}\\
2016--2023&Christopher Castro (Bioinformatics, University of Michigan, Chair)\\ 
&{\sl Investigating the Role of Noncoding De Novo Single-Nucleotide Variants in Autism Spectrum Disorder}\\
2017--2023&Melissa Englund (Genetics and Genomics, University of Michigan, Chair)\\ 
&{\sl Identification and Characterization of Cis-Regulatory Elements in the Human Genome}\\
2018--2023&Stephen Carney (Cancer Biology, University of Michigan, Committee Member)\\
&{\sl Epigenetic reprogramming in mutant IDH1 glioma influences radioresistance and neural lineage differentiation}\\
2019--2023&Benjamin Yang (Biomedical Engineering, University of Michigan, Committee Member)\\
&{\sl Towards Defining Principles of Cell Fate Plasticity}\\
2018--2022&Marcus Sherman (Bioinformatics, University of Michigan, Committee Member)\\
&{\sl Cultivation of enhanced bioinformatic-specific pedagogical manipulatives, interventions, and professional development}\\
2021--2022&Kuan-Han Hank Wu (Bioinformatics, University of Michigan, Committee Member)\\
&{\sl Integrating Electronic Health Records with Genetic Information to Advance Precision Medicine Approaches in Cardiovascular Disease}\\
2017--2022&Amanda Moccia (Genetics and Genomics, University of Michigan, Committee Member)\\
&{\sl Investigation of Developmental Disorders: Genetic Discovery and Functional Validation}\\
2017--2022&Ningxin Ouyang (Bioinformatics, University of Michigan, Chair)\\ 
&{\sl Deciphering Transcriptional Regulatory Circuits: Transcription Factor Binding and Regulatory Variants Identification}\\
2015--2021&Torrin McDonald (Genetics and Genomics, University of Michigan, Chair)\\ 
&{\sl Leveraging New Technologies to Explore Regulatory and Structural Elements of the Human Genome}\\
2018--2021&Heming Yao (Bioinformatics, University of Michigan, Committee Member)\\
&{\sl Machine Learning and Image Processing for Clinical Outcome Prediction: Applications in Medical Data from Patients with Traumatic Brain Injury, Ulcerative Colitis, and Heart Failure}\\
2016--2021&Mohd Hafiz Bin Mohd Rothi (Molecular, Cellular, and Developmental Biology, University of Michigan, Committee Member)\\
&{\sl Control of Chromatin by RNA-mediated Transcriptional Silencing}\\
2016--2021&Shengcheng Dong (Bioinformatics, University of Michigan, Chair)\\ 
&{\sl Computational Methods to Identify Regulatory Variants in the Non-coding Regions of the Human Genome}\\
2017--2021&Steven Romanelli (Molecular \& Integrative Physiology, University of Michigan, Committee Member)\\
&{\sl Viral CRISPR/Cas9 Gene Transfer for Somatic Knockout in Brown Adipose Tissue}\\
2018--2021&Negar Farzaneh (Bioinformatics, University of Michigan, Committee Member)\\
&{\sl Automated Decision Support System for Traumatic Injuries}\\
2016--2020&Shriya Sethuraman (Bioinformatics, University of Michigan, co-Chair)\\ 
&{\sl Genome-wide Identification of Non-coding Transcription by RNA Polymerase V and Its Involvement in Transcriptional Gene Silencing}\\
2015--2020&Sierra Nishizaki (Genetics and Genomics, University of Michigan, Chair)\\ 
&{\sl Decoding the Non-coding Genome: Novel Technologies for the Characterization of Non-coding Elements and Variation}\\
2017--2020&Christopher Lee (Biostatistics, University of Michigan, Committee Member)\\
&{\sl Improvements and Developments in Gene Regulation and Single-Cell Gene Expression Data Analysis}\\
2018--2019&Christine Ziegler (Biological Chemistry, University of Michigan, Committee Member)\\
2015--2018&Ari Allyn-Feuer (Bioinformatics, University of Michigan, Committee Member)\\
&{\sl The Pharmacoepigenomics Informatics Pipeline and H-GREEN Hi-C Compiler: Discovering Pharmacogenomic Variants and Pathways with the Epigenome and Spatial Genome}\\
2015--2017&Raymond Cavalcante (Bioinformatics, University of Michigan, Committee Member)\\
&{\sl Beyond the Transcriptome: Facilitating Interpretation of Epigenomics and Metabolomics Data}\\
2015--2017&Zhengting Zou (Bioinformatics, University of Michigan, Committee Member)\\
&{\sl Model-based genomic studies of protein sequence evolution: convergence, epistasis, and amino acid acceptance rates}\\
\end{longtable}

\subsubsection*{Preliminary Exam Committees (n=46)}
\begin{longtable}{L!{\VRule}R}
2026&James George (Cellular and Molecular Biology, University of Michigan)\\
2025&Xinyi Deng (Bioinformatics, University of Michigan)\\
2025&Mitchell Witt (Bioinformatics, University of Michigan)\\
2024&Luke Gohmann (Cellular and Molecular Biology, University of Michigan)\\
2024&Benjamin Li (Bioinformatics, University of Michigan)\\
2024&Tiffany Wan (Bioinformatics, University of Michigan)\\
2024&Zhiyuan Yu (Bioinformatics, University of Michigan)\\
2024&Rebecca McAvoy (Molecular, Cellular, and Developmental Biology, University of Michigan)\\
2024&Bonje Obua (Cellular and Molecular Biology, University of Michigan)\\
2024&Abigail Vallie (Cellular and Molecular Biology, University of Michigan)\\
2023&Jinhao Wang (Bioinformatics, University of Michigan)\\
2023&Lishi Yin (Bioinformatics, University of Michigan)\\
2023&Matthew Hodgman (Bioinformatics, University of Michigan)\\
2023&Ilakkiya Venkatachalam (Genetics and Genomics, University of Michigan)\\
2023&Jianhui Gong (Bioinformatics, University of Michigan)\\
2023&Mahnoor Gondal (Bioinformatics, University of Michigan)\\
2023&Elysia Chou (Bioinformatics, University of Michigan)\\
2022&Sean Moran (Bioinformatics, University of Michigan)\\
2022&Lu Lu (Bioinformatics, University of Michigan)\\
2022&Linghua Jiang (Bioinformatics, University of Michigan)\\
2022&Kaiwen Deng (Bioinformatics, University of Michigan)\\
2022&Yufeng Zhang (Bioinformatics, University of Michigan)\\
2021&Anthony Nguyen (Human Genetics, University of Michigan)\\
2021&Hanbyul Cho (Bioinformatics, University of Michigan)\\
2021&Charles Ryan (Cellular and Molecular Biology, University of Michigan)\\
2021&Kuan-Han Wu (Bioinformatics, University of Michigan)\\
2021&Wenjin Gu (Bioinformatics, University of Michigan)\\
2020&Jie Cao (Bioinformatics, University of Michigan)\\
2020&Zijun Gao (Bioinformatics, University of Michigan)\\
2020&Ashley Melnick (Cellular and Molecular Biology, University of Michigan)\\
2019&Benjamin Yang (Biomedical Engineering, University of Michigan)\\
2019&Maria Virgilio (Cellular and Molecular Biology, University of Michigan)\\
2018&Zhi Carrie Li (Bioinformatics, University of Michigan)\\
2018&Kevin Hu (Bioinformatics, University of Michigan)\\
2018&Siyu Liu (Bioinformatics, University of Michigan)\\
2018&Alexandra Weber (Bioinformatics, University of Michigan)\\
2018&Mitch Fernandez (Bioinformatics, University of Michigan)\\
2017&Tingyang Li (Bioinformatics, University of Michigan)\\
2017&Marcus Sherman (Bioinformatics, University of Michigan)\\
2017&Adrienne Shami (Human Genetics, University of Michigan)\\
2017&Trenton Frisbie (Human Genetics, University of Michigan)\\
2017&Melissa Englund (Human Genetics, University of Michigan)\\
2017&Peter Orchard (Bioinformatics, University of Michigan)\\
2017&Li Guan (Bioinformatics, University of Michigan)\\
2016&Shriya Sethuraman (Bioinformatics, University of Michigan)\\
2016&Jed Carlson (Bioinformatics, University of Michigan)\\
\end{longtable}

\section*{Industry Experience}
\begin{longtable}{L!{\VRule}R}
2013--2014&Consultant, Color Genomics\\
&Personalized medicine / genomics startup\\
\end{longtable}

% Publications and patents are generated from the canonical website data.
% Generated by scripts/build_cv.py. Do not edit directly.
% Source: bibliography/publications.bib + publication_metadata + _people.
\cvsection{Publications}
\begin{flushright}
\scriptsize * co-first authorship; \(\dagger\) co-senior authorship; \underline{underline} indicates Boyle Lab members
\end{flushright}
\begin{enumerate}[leftmargin=2.7em,labelsep=0.55em,itemsep=0.65em,parsep=0pt,topsep=0.25em]
\item[{[87]}] \underline{Pavlovic K}, \underline{McDonald TL}, \underline{Diehl AG}, \underline{Switzenberg JA}, \textbf{Boyle AP}. ``Fiber-TEnCATS reveals haplotype-specific chromatin accessibility and DNA methylation at human L1HS loci.'' \textit{bioRxiv} 2026. \href{\detokenize{https://doi.org/10.64898/2026.06.26.734832}}{doi: 10.64898/2026.06.26.734832}.
\item[{[86]}] \underline{Diehl AG}, \textbf{Boyle AP}. ``GPatch enables chromosome-scale, gap-free pseudoassemblies from fragmented draft genomes.'' \textit{bioRxiv} 2026. \href{\detokenize{https://doi.org/10.1101/2025.05.22.655567}}{doi: 10.1101/2025.05.22.655567}.
\item[{[85]}] Kenney GE, \underline{Sherpa RN}, Burgess JD, \textbf{Boyle AP}, Phanstiel DH. ``SEMPLR: an R package for transcription factor binding prediction.'' \textit{Bioinformatics} 2026, 42(6):btag383. \href{\detokenize{https://pubmed.ncbi.nlm.nih.gov/42286344/}}{PMID: 42286344}.
\item[{[84]}] Rabbani M, Apell Z, Parnell TJ, Moritz L, Kim S, \underline{Srinivasan S}, Agrawal R, Vargo A, Orchard P, Xie W, Freddolino L, \textbf{Boyle AP}, Li JZ, Lesch BJ, Cairns B, Kim M, Wilson TE, Hammoud SS. ``3D chromatin compartment of round spermatids encodes the spatiotemporal program of histone-to-protamine exchange in spermiogenesis.'' \textit{bioRxiv} 2026. \href{\detokenize{https://doi.org/10.64898/2026.03.10.710708}}{doi: 10.64898/2026.03.10.710708}.
\item[{[83]}] Li S, Chou E, Wang K, \textbf{Boyle AP}, Sartor MA. ``TFBSpedia: a comprehensive human and mouse transcription factor binding sites database.'' \textit{bioRxiv} 2026. \href{\detokenize{https://doi.org/10.64898/2026.03.04.709638}}{doi: 10.64898/2026.03.04.709638}.
\item[{[82]}] \underline{McBean B}, Hauk B, Michmerhuizen AR, Chandler BC, Pesch AM, Lerner LM, Gurdak D, Ward C, Rana P, Zeidane RA, Mercer V, Tao M, Hochmuth E, Jungles KM, The S, Liu M, Spratt DE, \textbf{Boyle AP}, Pierce LJ, Speers CW. ``Transcriptomic analysis to uncover the mechanism of radiosensitization of AR-positive triple-negative breast cancers with AR inhibition.'' \textit{NPJ Breast Cancer} 2026, 12(1). \href{\detokenize{https://pubmed.ncbi.nlm.nih.gov/41735341/}}{PMID: 41735341}.
\item[{[81]}] Li D, Liu S, Assis PR, Li M, Dong S, Whaling I, Jolanki O, Kagda M, Zhang W, Macias-Velasco JF, Liu T, Cody S, Antonacci-Fulton L, Huang Y, Liu J, Montgomery MT, Zeiberg D, Jain S, Pejaver V, Bergquist T, Chen Y, Radivojac P, Gersbach CA, \underline{Sherpa RN}, \underline{Castro CP}, \textbf{Boyle AP}, Starita LM, Fowler DM, Ahituv N, Dey KK, Majoros WH, Reddy TE, Craven M, Sinha R, Sverchkov Y, Cai X, Nzima MZ, Calderwood MA, Rozowsky J, Gerstein M, Ma J, Yue F, Cherry JM, Love MI, Engreitz JM, Hitz BC, Wang T. ``The IGVF catalog—from genetic variation to function.'' \textit{Nucleic Acids Research} 2025, gkaf1341. \href{\detokenize{https://pubmed.ncbi.nlm.nih.gov/41359121/}}{PMID: 41359121}.
\item[{[80]}] The Somatic Mosaicism across Human Tissues Network (SMaHT). ``Comprehensive benchmarking of somatic mutation detection by the SMaHT Network.'' \textit{bioRxiv} 2025. \href{\detokenize{https://doi.org/10.1101/2025.10.09.678885}}{doi: 10.1101/2025.10.09.678885}.
\item[{[79]}] *Wang S, *Bae M, *Wang J, Zhao B, Nguyen K, Mallett S, \underline{Switzenberg JA}, \underline{Losh SJ}, Sexton CE, Miao B, Dong S, Zeng X, Wang Z, \underline{McDonald TL}, \underline{Mumm C}, Gadde RK, Tariq AM, Chen Z, Feng WC, Burn A, Park J, Chu C, Shen H, Wang T, Urban AE, Zhu X, Li H, Burns KH, Chun HJE, Park PJ, SMaHT MEI Working Group, \textsuperscript{\(\dagger\)}\textbf{Boyle AP}, \textsuperscript{\(\dagger\)}Mills RE, \textsuperscript{\(\dagger\)}Zhou W, \textsuperscript{\(\dagger\)}Lee EA. ``Multi-platform framework for mapping somatic retrotransposition in‬ human tissues.'' \textit{bioRxiv} 2025. \href{\detokenize{https://doi.org/10.1101/2025.10.07.680917}}{doi: 10.1101/2025.10.07.680917}.
\item[{[78]}] Denstaedt SJ, \underline{McBean B}, \textbf{Boyle AP}, Arenberg BC, Mack M, Moore BB, Newstead MW, Deng Y, Nesvizhskii AI, Singer BH, Cano J, Prescott HC, Goodridge HS, Zemans RL. ``Long-term immune reprogramming of classical monocytes with altered ontogeny mediates enhanced lung injury in sepsis survivors.'' \textit{bioRxiv} 2025. \href{\detokenize{https://doi.org/10.1101/2025.05.16.654442}}{doi: 10.1101/2025.05.16.654442}.
\item[{[77]}] Coorens THH, Oh JW, Choi YA, Lim NS, Zhao B, Voshall A, Abyzov A, Antonacci-Fulton L, Aparicio S, Ardlie KG, Bell TJ, Bennett JT, Bernstein BE, Blanchard TG, \textbf{Boyle AP}, Buenrostro JD, Burns KH, Chen F, Chen R, Choudhury S, Doddapaneni HV, Eichler EE, Evrony GD, Faith MA, Fazzio TG, Fulton RS, Garber M, Gehlenborg N, Germer S, Getz G, Gibbs RA, Hernandez RG, Jin F, Korbel JO, Landau DA, Lawson HA, Lennon NJ, Li H, Li Y, Loh PR, Marth G, McConnell MJ, Mills RE, Montgomery SB, Natarajan P, Park PJ, Satija R, Sedlazeck FJ, Shao DD, Shen H, Stergachis AB, Underhill HR, Urban AE, VonDran MW, Walsh CA, Wang T, Wu TP, Zong C, Lee EA, Vaccarino FM, Somatic Mosaicism across Human Tissues Network. ``The Somatic Mosaicism across Human Tissues Network.'' \textit{Nature} 2025, 643(8070):47-59. \href{\detokenize{https://pubmed.ncbi.nlm.nih.gov/40604182/}}{PMID: 40604182}.
\item[{[76]}] *\underline{Van Deynze K}, *\underline{Mumm C}, Maltby CJ, \underline{Switzenberg JA}, Todd PK, \textbf{Boyle AP}. ``Enhanced detection and genotyping of disease-associated tandem repeats using HMMSTR and targeted long-read sequencing.'' \textit{Nucleic Acids Research} 2025, 53(2):gkae1202. \href{\detokenize{https://pubmed.ncbi.nlm.nih.gov/39676678/}}{PMID: 39676678}.
\item[{[75]}] \underline{Parana P}, \underline{Mumm C}, McConnell MJ, \textbf{Boyle AP}. ``Draft de novo genome construction of Scytonema sp. PRP1: identified from single-cell sequencing library preparation.'' \textit{Microbiology Resource Announcements} 2025, e00029-25. \href{\detokenize{https://pubmed.ncbi.nlm.nih.gov/40744872/}}{PMID: 40744872}.
\item[{[74]}] Miller SL, Green KM, \underline{Crone B}, \underline{Switzenberg JA}, Tank EMH, Krans A, Jansen-West K, Wieland CM, Ji EW, Petrucelli L, Barmada SJ, \textbf{Boyle AP}, Todd PK. ``Cryptic intronic transcriptional initiation generates efficient endogenous mRNA templates for C9orf72-associated RAN translation.'' \textit{Proceedings of the National Academy of Sciences} 2025, 122(32):e2507334122. \href{\detokenize{https://pubmed.ncbi.nlm.nih.gov/40758885/}}{PMID: 40758885}.
\item[{[73]}] \underline{McBean B}, Abou Zeidane R, Lichtman-Mikol S, Hauk B, Speers J, Tidmore S, Flores CL, Rana PS, Pisano C, Liu M, Santola A, Montero A, \textbf{Boyle AP}, Speers CW. ``MELK as a Mediator of Stemness and Metastasis in Aggressive Subtypes of Breast Cancer.'' \textit{International Journal of Molecular Sciences} 2025, 26(5):2245. \href{\detokenize{https://pubmed.ncbi.nlm.nih.gov/40076867/}}{PMID: 40076867}.
\item[{[72]}] Lee S, McAfee JC, Lee J, Gomez A, Ledford AT, Clarke D, Min H, Gerstein MB, \textbf{Boyle AP}, Sullivan PF, Beltran AS, Won H. ``Massively parallel reporter assay investigates shared genetic variants of eight psychiatric disorders.'' \textit{Cell} 2025, S0092-8674(24):01435-1. \href{\detokenize{https://pubmed.ncbi.nlm.nih.gov/39848247/}}{PMID: 39848247}.
\item[{[71]}] *Zhou W, *\underline{Mumm C}, *Gan Y, \underline{Switzenberg JA}, Wang J, De Oliveira P, Kathuria K, \underline{Losh SJ}, \underline{McDonald TL}, Bessell B, \underline{Van Deynze K}, \textsuperscript{\(\dagger\)}McConnell MJ, \textsuperscript{\(\dagger\)}\textbf{Boyle AP}, \textsuperscript{\(\dagger\)}Mills RE. ``A personalized multi-platform assessment of somatic mosaicism in the human frontal cortex.'' \textit{bioRxiv} 2024. \href{\detokenize{https://doi.org/10.1101/2024.12.18.629274}}{doi: 10.1101/2024.12.18.629274}.
\item[{[70]}] Maltby CJ, Krans A, Grudzien SJ, Palacios Y, Muiños J, Suárez A, Asher M, Willey S, \underline{Van Deynze K}, \underline{Mumm C}, \textbf{Boyle AP}, Cortese A, Khurana V, Barmada SJ, Dijkstra AA, Todd PK. ``AAGGG repeat expansions trigger RFC1-independent synaptic dysregulation in human CANVAS neurons.'' \textit{Science Advances} 2024, 10(36):eadn2321. \href{\detokenize{https://pubmed.ncbi.nlm.nih.gov/39231235/}}{PMID: 39231235}.
\item[{[69]}] IGVF Consortium. ``Deciphering the impact of genomic variation on function.'' \textit{Nature} 2024, 633:47--57. \href{\detokenize{https://pubmed.ncbi.nlm.nih.gov/39232149/}}{PMID: 39232149}.
\item[{[68]}] Yee C, Xiao Y, Chen H, Reddy A, Xu B, Medwig-Kinney T, Zhang W, \textbf{Boyle AP}, Herbst W, Xiang Y, Matus D, Shen K. ``An activity-regulated transcriptional program directly drives synaptogenesis.'' \textit{Nature Neuroscience} 2024, 27(9):1695-1707. \href{\detokenize{https://pubmed.ncbi.nlm.nih.gov/39103556/}}{PMID: 39103556}.
\item[{[67]}] \underline{Crone B}, \textbf{Boyle AP}. ``Enhancing Portability of Trans-Ancestral Polygenic Risk Scores through Tissue-Specific Functional Genomic Data Integration.'' \textit{PLoS Genetics} 2024, 20(8):e1011356. \href{\detokenize{https://pubmed.ncbi.nlm.nih.gov/39110742/}}{PMID: 39110742}.
\item[{[66]}] \underline{Holmes MJ}, Mahjour B, \underline{Castro CP}, \underline{Farnum GA}, \underline{Diehl AG}, \textbf{Boyle AP}. ``HaplotagLR: an efficient and configurable utility for haplotagging long reads.'' \textit{PLoS ONE} 2024, 19(3):1-15. \href{\detokenize{https://pubmed.ncbi.nlm.nih.gov/38478504/}}{PMID: 38478504}.
\item[{[65]}] The Critical Assessment of Genome Interpretation Consortium. ``CAGI, the Critical Assessment of Genome Interpretation, establishes progress and prospects for computational genetic variant interpretation methods.'' \textit{Genome Biology} 2024, 25(1):53. \href{\detokenize{https://pubmed.ncbi.nlm.nih.gov/38389099/}}{PMID: 38389099}.
\item[{[64]}] McAfee JC, Lee S, Lee J, Bell JL, Krupa O, Davis J, Insigne K, Bond ML, \underline{Zhao N}, \textbf{Boyle AP}, Phanstiel DH, Love MI, Stein JL, Ruzicka WB, Davila-Velderrain J, Kosuri S, Won H. ``Systematic investigation of allelic regulatory activity of schizophrenia-associated common variants.'' \textit{Cell Genomics} 2023, 3:100404. \href{\detokenize{https://pubmed.ncbi.nlm.nih.gov/37868037/}}{PMID: 37868037}.
\item[{[63]}] \underline{Zhao N}, \underline{Dong S}, \textbf{Boyle AP}. ``Organ-specific prioritization and annotation of non-coding regulatory variants in the human genome.'' \textit{bioRxiv} 2023. \href{\detokenize{https://doi.org/10.1101/2023.09.07.556700}}{doi: 10.1101/2023.09.07.556700}.
\item[{[62]}] Moritz L, Schon SB, Rabbani M, Sheng Y, Agrawal R, Glass-Klaiber J, Sultan C, M CJ, Clements J, Baldwin MR, \underline{Diehl AG}, \textbf{Boyle AP}, O'Brien PJ, Ragunathan K, Hu YC, Kelleher NL, Nandakumar J, Li JZ, Orwig KE, Redding S, Hammoud SS. ``Sperm chromatin structure and reproductive fitness are altered by substitution of a single amino acid in mouse protamine 1.'' \textit{Nature Structural \& Molecular Biology} 2023. \href{\detokenize{https://pubmed.ncbi.nlm.nih.gov/37460896/}}{PMID: 37460896}.
\item[{[61]}] *\underline{Mumm C}, *\underline{Drexel ML}, \underline{McDonald TL}, \underline{Diehl AG}, \underline{Switzenberg JA}, \textbf{Boyle AP}. ``OnRamp: rapid nanopore plasmid validation.'' \textit{Genome Research} 2023, 33(5):741--749. \href{\detokenize{https://pubmed.ncbi.nlm.nih.gov/37156622/}}{PMID: 37156622}.
\item[{[60]}] \underline{Castro CP}, \underline{Diehl AG}, \textbf{Boyle AP}. ``Challenges in screening for de novo noncoding variants contributing to genetically complex phenotypes.'' \textit{Human Genetics and Genomics Advances} 2023, 4(3):100210. \href{\detokenize{https://pubmed.ncbi.nlm.nih.gov/37305558/}}{PMID: 37305558}.
\item[{[59]}] *\underline{Dong S}, *\underline{Zhao N}, Spragins E, Kagda MS, Li M, Assis PR, Jolanki O, Luo Y, Cherry JM, \textsuperscript{\(\dagger\)}\textbf{Boyle AP}, \textsuperscript{\(\dagger\)}Hitz BC. ``Annotating and prioritizing human non-coding variants with RegulomeDB v.2.'' \textit{Nature Genetics} 2023, 55(5):724--726. \href{\detokenize{https://pubmed.ncbi.nlm.nih.gov/37173523/}}{PMID: 37173523}.
\item[{[58]}] \underline{Zhao N}, Wang S, Huang Q, \underline{Dong S}, \textbf{Boyle AP}. ``Explain-seq: an end-to-end pipeline from training to interpretation of sequence-based deep learning models.'' \textit{bioRxiv} 2023. \href{\detokenize{https://doi.org/10.1101/2023.01.23.525250}}{doi: 10.1101/2023.01.23.525250}.
\item[{[57]}] \underline{Ouyang N}, \textbf{Boyle AP}. ``Quantitative assessment of association between noncoding variants and transcription factor binding.'' \textit{bioRxiv} 2022. \href{\detokenize{https://doi.org/10.1101/2022.11.22.517559}}{doi: 10.1101/2022.11.22.517559}.
\item[{[56]}] \underline{Nishizaki SS}, \textbf{Boyle AP}. ``SEMplMe: A tool for integrating DNA methylation effects in transcription factor binding affinity predictions.'' \textit{BMC Bioinformatics} 2022, 23(1):317. \href{\detokenize{https://pubmed.ncbi.nlm.nih.gov/35927613/}}{PMID: 35927613}.
\item[{[55]}] Qin T, Lee C, Li S, Cavalcante RG, Orchard P, Yao H, Zhang H, Wang S, Patil S, \textbf{Boyle AP}, Sartor MA. ``Comprehensive enhancer-target gene assignments improve gene set level interpretation of genome-wide regulatory data.'' \textit{Genome Biology} 2022, 23(1):105. \href{\detokenize{https://pubmed.ncbi.nlm.nih.gov/35473573/}}{PMID: 35473573}.
\item[{[54]}] \underline{Dong S}, \textbf{Boyle AP}. ``Prioritization of regulatory variants with tissue-specific function in the non-coding regions of human genome.'' \textit{Nucleic Acids Research} 2021, 50(1):e6-e6. \href{\detokenize{https://pubmed.ncbi.nlm.nih.gov/34648033/}}{PMID: 34648033}.
\item[{[53]}] Bao Y, Wadden J, Erb-Downward JR, Ranjan P, Zhou W, \underline{McDonald TL}, Mills RE, \textbf{Boyle AP}, Dickson RP, Blaauw D, Welch JD. ``SquiggleNet: real-time, direct classification of nanopore signals.'' \textit{Genome Biology} 2021, 22(1):298. \href{\detokenize{https://pubmed.ncbi.nlm.nih.gov/34706748/}}{PMID: 34706748}.
\item[{[52]}] *\underline{McDonald TL}, *Zhou W, \underline{Castro CP}, \underline{Mumm C}, \underline{Switzenberg JA}, \textsuperscript{\(\dagger\)}Mills RE, \textsuperscript{\(\dagger\)}\textbf{Boyle AP}. ``Cas9 targeted enrichment of mobile elements using nanopore sequencing.'' \textit{Nature Communications} 2021, 12(1):3586. \href{\detokenize{https://pubmed.ncbi.nlm.nih.gov/34117247/}}{PMID: 34117247}.
\item[{[51]}] \underline{Zhao N}, \textbf{Boyle AP}. ``F-Seq2: improving the feature density based peak caller with dynamic statistics.'' \textit{NAR Genomics and Bioinformatics} 2021, 3(1). \href{\detokenize{https://pubmed.ncbi.nlm.nih.gov/33655209/}}{PMID: 33655209}.
\item[{[50]}] *\underline{Nishizaki SS}, *\underline{McDonald TL}, \underline{Farnum GA}, \underline{Holmes MJ}, \underline{Drexel ML}, \underline{Switzenberg JA}, \textbf{Boyle AP}. ``The inducible lac operator-repressor system is functional in zebrafish cells.'' \textit{Frontiers in Genetics} 2021, 12. \href{\detokenize{https://pubmed.ncbi.nlm.nih.gov/34220959/}}{PMID: 34220959}.
\item[{[49]}] *Tsuzuki M, *\underline{Sethuraman S}, Coke AN, Rothi MH, \textbf{Boyle AP}, Wierzbicki AT. ``Broad noncoding transcription suggests genome surveillance by RNA polymerase V.'' \textit{Proceedings of the National Academy of Sciences} 2020, 117(48):30799--30804. \href{\detokenize{https://pubmed.ncbi.nlm.nih.gov/33199612/}}{PMID: 33199612}.
\item[{[48]}] Rothi MH, \underline{Sethuraman S}, Dolata J, \textbf{Boyle AP}, Wierzbicki AT. ``DNA methylation directs nucleosome positioning in RNA-mediated transcriptional silencing.'' \textit{bioRxiv} 2020. \href{\detokenize{https://doi.org/10.1101/2020.10.29.359794}}{doi: 10.1101/2020.10.29.359794}.
\item[{[47]}] \underline{Diehl AG}, \textbf{Boyle AP}. ``MapGL: Inferring evolutionary gain and loss of short genomic sequence features by phylogenetic maximum parsimony.'' \textit{BMC Bioinformatics} 2020, 21:416. \href{\detokenize{https://pubmed.ncbi.nlm.nih.gov/32962625/}}{PMID: 32962625}.
\item[{[46]}] \underline{Ouyang N}, \textbf{Boyle AP}. ``TRACE: transcription factor footprinting using chromatin accessibility data and DNA sequence.'' \textit{Genome Research} 2020, 30(7):1040--1046. \href{\detokenize{https://pubmed.ncbi.nlm.nih.gov/32660981/}}{PMID: 32660981}.
\item[{[45]}] The ENCODE Project Consortium. ``Perspectives on ENCODE.'' \textit{Nature} 2020, 583(7818):693--698. \href{\detokenize{https://pubmed.ncbi.nlm.nih.gov/32728248/}}{PMID: 32728248}.
\item[{[44]}] The ENCODE Project Consortium. ``Expanded encyclopaedias of DNA elements in the human and mouse genomes.'' \textit{Nature} 2020, 583(7818):699--710. \href{\detokenize{https://pubmed.ncbi.nlm.nih.gov/32728249/}}{PMID: 32728249}.
\item[{[43]}] \underline{Diehl AG}, \underline{Ouyang N}, \textbf{Boyle AP}. ``Transposable elements contribute to cell and species-specific chromatin looping and gene regulation in mammalian genomes.'' \textit{Nature Communications} 2020, 11(1):1796. \href{\detokenize{https://pubmed.ncbi.nlm.nih.gov/32286261/}}{PMID: 32286261}.
\item[{[42]}] Lee CT, Cavalcante RG, Lee C, Qin T, Patil S, Wang S, Tsai ZTY, \textbf{Boyle AP}, Sartor MA. ``Poly-Enrich: count-based methods for gene set enrichment testing with genomic regions.'' \textit{NAR Genomics and Bioinformatics} 2020, 2(1). \href{\detokenize{https://pubmed.ncbi.nlm.nih.gov/32051932/}}{PMID: 32051932}.
\item[{[41]}] \underline{Nishizaki SS}, Ng N, \underline{Dong S}, Porter RS, \underline{Morterud C}, \underline{Williams C}, \underline{Asman C}, \underline{Switzenberg JA}, \textbf{Boyle AP}. ``Predicting the effects of SNPs on transcription factor binding affinity.'' \textit{Bioinformatics} 2019, 50:2434. \href{\detokenize{https://pubmed.ncbi.nlm.nih.gov/31373606/}}{PMID: 31373606}.
\item[{[40]}] \underline{Diehl AG}, \textbf{Boyle AP}. ``CGIMP: Real-time exploration and covariate projection for self-organizing map datasets.'' \textit{Journal of Open Source Software} 2019, 4(39):1520. \href{\detokenize{https://doi.org/10.21105/joss.01520}}{doi: 10.21105/joss.01520}.
\item[{[39]}] \underline{Amemiya HM}, \textsuperscript{\(\dagger\)}Kundaje A, \textsuperscript{\(\dagger\)}\textbf{Boyle AP}. ``The ENCODE Blacklist: Identification of Problematic Regions of the Genome.'' \textit{Scientific Reports} 2019, 9(1):9354. \href{\detokenize{https://pubmed.ncbi.nlm.nih.gov/31249361/}}{PMID: 31249361}.
\item[{[38]}] Varshney A, VanRenterghem H, Orchard P, \textsuperscript{\(\dagger\)}\textbf{Boyle AP}, \textsuperscript{\(\dagger\)}Stitzel ML, \textsuperscript{\(\dagger\)}Ucar D, Parker SC. ``Cell specificity of regulatory annotations and their genetic effects on gene expression.'' \textit{Genetics} 2019, 211(2):549--562. \href{\detokenize{https://pubmed.ncbi.nlm.nih.gov/30593493/}}{PMID: 30593493}.
\item[{[37]}] Shigaki D, Adato O, Adhikar AN, \underline{Dong S}, Hawkins-Hooker A, Inoue F, Juven-Gershon T, Kenlay H, Martin B, Patra A, Penzar DP, Schubach M, Xiong C, Yan Z, \textbf{Boyle AP}, Kreimer A, Kulakovskiy IV, Reid J, Unger R, Yosef N, Shendure J, Ahituv N, Kircher M, Beer MA. ``Integration of Multiple Epigenomic Marks Improves Prediction of Variant Impact in Saturation Mutagenesis Reporter Assay.'' \textit{Human mutation} 2019, 33(8):831. \href{\detokenize{https://pubmed.ncbi.nlm.nih.gov/31106481/}}{PMID: 31106481}.
\item[{[36]}] \underline{Dong S}, \textbf{Boyle AP}. ``Predicting functional variants in enhancer and promoter elements using RegulomeDB.'' \textit{Human Mutation} 2019, 33(8):831. \href{\detokenize{https://pubmed.ncbi.nlm.nih.gov/31228310/}}{PMID: 31228310}.
\item[{[35]}] \underline{Diehl AG}, \textbf{Boyle AP}. ``Conserved and species-specific transcription factor co-binding patterns drive divergent gene regulation in human and mouse.'' \textit{Nucleic Acids Research} 2018, 46(4):1878--1894. \href{\detokenize{https://pubmed.ncbi.nlm.nih.gov/29361190/}}{PMID: 29361190}.
\item[{[34]}] Nielsen JB, Fritsche LG, Zhou W, Teslovich TM, Holmen OL, Gustafsson S, Gabrielsen ME, Schmidt EM, Beaumont R, Wolford BN, Lin M, Brummett CM, Preuss MH, Refsgaard L, Bottinger EP, Graham SE, Surakka I, Chu Y, Skogholt AH, Dalen H, \textbf{Boyle AP}, Oral H, Herron TJ, Kitzman J, Jalife J, Svendsen JH, Olesen MS, Njølstad I, Løchen ML, Baras A, Gottesman O, Marcketta A, O'Dushlaine C, Ritchie MD, Wilsgaard T, Loos RJF, Frayling TM, Boehnke M, Ingelsson E, Carey DJ, Dewey FE, Kang HM, Abecasis GR, Hveem K, Willer CJ. ``Genome-wide Study of Atrial Fibrillation Identifies Seven Risk Loci and Highlights Biological Pathways and Regulatory Elements Involved in Cardiac Development.'' \textit{American Journal of Human Genetics} 2017, 102(1):103--115. \href{\detokenize{https://pubmed.ncbi.nlm.nih.gov/29290336/}}{PMID: 29290336}.
\item[{[33]}] Spadafore M, Najarian K, \textbf{Boyle AP}. ``A proximity-based graph clustering method for the identification and application of transcription factor clusters.'' \textit{BMC Bioinformatics} 2017, 18(1):530. \href{\detokenize{https://pubmed.ncbi.nlm.nih.gov/29187152/}}{PMID: 29187152}.
\item[{[32]}] *Yang B, *Zhou W, *Jiao J, Nielsen JB, Mathis MR, Heydarpour M, Lettre G, Folkersen L, Prakash S, Schurmann C, Fritsche L, \underline{Farnum GA}, Lin M, Othman M, Hornsby W, Driscoll A, Levasseur A, Thomas M, Farhat L, Dubé MP, Isselbacher EM, Franco-Cereceda A, Guo DC, Bottinger EP, Deeb GM, Booher A, Kheterpal S, Chen YE, Kang HM, Kitzman J, Cordell HJ, Keavney BD, Goodship JA, Ganesh SK, Abecasis G, Eagle KA, \textbf{Boyle AP}, Loos RJF, \textsuperscript{\(\dagger\)}Eriksson P, \textsuperscript{\(\dagger\)}Tardif JC, \textsuperscript{\(\dagger\)}Brummett CM, \textsuperscript{\(\dagger\)}Milewicz DM, \textsuperscript{\(\dagger\)}Body SC, \textsuperscript{\(\dagger\)}Willer CJ. ``Protein-altering and regulatory genetic variants near GATA4 implicated in bicuspid aortic valve.'' \textit{Nature Communications} 2017, 8:15481. \href{\detokenize{https://pubmed.ncbi.nlm.nih.gov/28541271/}}{PMID: 28541271}.
\item[{[31]}] \underline{Nishizaki SS}, \textbf{Boyle AP}. ``Mining the Unknown: Assigning Function to Noncoding Single Nucleotide Polymorphisms.'' \textit{Trends in Genetics} 2017, 33(1):34--45. \href{\detokenize{https://pubmed.ncbi.nlm.nih.gov/27939749/}}{PMID: 27939749}.
\item[{[30]}] \underline{Diehl AG}, \textbf{Boyle AP}. ``Deciphering ENCODE.'' \textit{Trends in Genetics} 2016, 32(4):238--249. \href{\detokenize{https://pubmed.ncbi.nlm.nih.gov/26962025/}}{PMID: 26962025}.
\item[{[29]}] Phanstiel DH, \textbf{Boyle AP}, Heidari N, Snyder MP. ``Mango: A bias correcting ChIA-PET analysis pipeline.'' \textit{Bioinformatics} 2015. \href{\detokenize{https://pubmed.ncbi.nlm.nih.gov/26034063/}}{PMID: 26034063}.
\item[{[28]}] *Yue F, *Cheng Y, *Breschi A, *Vierstra J, *Wu W, *Ryba T, *Sandstrom R, *Ma Z, *Davis C, *Pope BD, *Shen Y, Pervouchine DD, Djebali S, Thurman RE, Kaul R, Rynes E, Kirilusha A, Marinov GK, Williams BA, Trout D, Amrhein H, Fisher-Aylor K, Antoshechkin I, DeSalvo G, See LH, Fastuca M, Drenkow J, Zaleski C, Dobin A, Prieto P, Lagarde J, Bussotti G, Tanzer A, Denas O, Li K, Bender MA, Zhang M, Byron R, Groudine MT, McCleary D, Pham L, Ye Z, Kuan S, Edsall L, Wu YC, Rasmussen MD, Bansal MS, Kellis M, Keller CA, Morrissey CS, Mishra T, Jain D, Dogan N, Harris RS, Cayting P, Kawli T, \textbf{Boyle AP}, Euskirchen G, Kundaje A, Lin S, Lin Y, Jansen C, Malladi VS, Cline MS, Erickson DT, Kirkup VM, Learned K, Sloan CA, Rosenbloom KR, Lacerda de Sousa B, Beal K, Pignatelli M, Flicek P, Lian J, Kahveci T, Lee D, Kent WJ, Ramalho Santos M, Herrero J, Notredame C, Johnson A, Vong S, Lee K, Bates D, Neri F, Diegel M, Canfield T, Sabo PJ, Wilken MS, Reh TA, Giste E, Shafer A, Kutyavin T, Haugen E, Dunn D, Reynolds AP, Neph S, Humbert R, Hansen RS, De Bruijn M, Selleri L, Rudensky A, Josefowicz S, Samstein R, Eichler EE, Orkin SH, Levasseur D, Papayannopoulou T, Chang KH, Skoultchi A, Gosh S, Disteche C, Treuting P, Wang Y, Weiss MJ, Blobel GA, Cao X, Zhong S, Wang T, Good PJ, Lowdon RF, Adams LB, Zhou XQ, Pazin MJ, Feingold EA, Wold B, Taylor J, Mortazavi A, Weissman SM, Stamatoyannopoulos JA, Snyder MP, Guigo R, Gingeras TR, Gilbert DM, Hardison RC, Beer MA, Ren B, Mouse ENCODE Consortium. ``A comparative encyclopedia of DNA elements in the mouse genome.'' \textit{Nature} 2014, 515(7527):355--364. \href{\detokenize{https://pubmed.ncbi.nlm.nih.gov/25409824/}}{PMID: 25409824}.
\item[{[27]}] *Cheng Y, *Ma Z, Kim BH, Wu W, Cayting P, \textbf{Boyle AP}, Sundaram V, Xing X, Dogan N, Li J, Euskirchen G, Lin S, Lin Y, Visel A, Kawli T, Yang X, Patacsil D, Keller CA, Giardine B, Mouse ENCODE Consortium, Kundaje A, Wang T, Pennacchio LA, Weng Z, \textsuperscript{\(\dagger\)}Hardison RC, \textsuperscript{\(\dagger\)}Snyder MP. ``Principles of regulatory information conservation between mouse and human.'' \textit{Nature} 2014, 515(7527):371--375. \href{\detokenize{https://pubmed.ncbi.nlm.nih.gov/25409826/}}{PMID: 25409826}.
\item[{[26]}] *\textbf{Boyle AP}, *Araya CL, Brdlik C, Cayting P, Cheng C, Cheng Y, Gardner K, Hillier LW, Janette J, Jiang L, Kasper D, Kawli T, Kheradpour P, Kundaje A, Li JJ, Ma L, Niu W, Rehm EJ, Rozowsky J, Slattery M, Spokony R, Terrell R, Vafeados D, Wang D, Weisdepp P, Wu YC, Xie D, Yan KK, Feingold EA, Good PJ, Pazin MJ, Huang H, Bickel PJ, Brenner SE, Reinke V, Waterston RH, Gerstein M, \textsuperscript{\(\dagger\)}White KP, \textsuperscript{\(\dagger\)}Kellis M, \textsuperscript{\(\dagger\)}Snyder M. ``Comparative analysis of regulatory information and circuits across distant species.'' \textit{Nature} 2014, 512(7515):453--456. \href{\detokenize{https://pubmed.ncbi.nlm.nih.gov/25164757/}}{PMID: 25164757}.
\item[{[25]}] Araya CL, Kawli T, Kundaje A, Jiang L, Wu B, Vafeados D, Terrell R, Weissdepp P, Gevirtzman L, Mace D, Niu W, \textbf{Boyle AP}, Xie D, Ma L, Murray JI, Reinke V, Waterston RH, Snyder M. ``Regulatory analysis of the C. elegans genome with spatiotemporal resolution.'' \textit{Nature} 2014, 512(7515):400--405. \href{\detokenize{https://pubmed.ncbi.nlm.nih.gov/25164749/}}{PMID: 25164749}.
\item[{[24]}] Phanstiel DH, \textbf{Boyle AP}, Araya CL, Snyder MP. ``Sushi.R: flexible, quantitative and integrative genomic visualizations for publication-quality multi-panel figures.'' \textit{Bioinformatics} 2014. \href{\detokenize{https://pubmed.ncbi.nlm.nih.gov/24903420/}}{PMID: 24903420}.
\item[{[23]}] *Kasowski M, *Kyriazopoulou-Panagiotopoulou S, *Grubert F, *Zaugg JB, *Kundaje A, Liu Y, \textbf{Boyle AP}, Zhang QC, Zakharia F, Spacek DV, Li J, Xie D, Steinmetz LM, Hogenesch JB, Kellis M, Batzoglou S, Snyder M. ``Extensive variation in chromatin states across humans.'' \textit{Science} 2013, 342(6159):750--752. \href{\detokenize{https://pubmed.ncbi.nlm.nih.gov/24136358/}}{PMID: 24136358}.
\item[{[22]}] *Xie D, *\textbf{Boyle AP}, *Wu L, Kawli T, Zhai J, Snyder M. ``Dynamic trans-acting factor colocalization in human cells.'' \textit{Cell} 2013, 155(3):713--724. \href{\detokenize{https://pubmed.ncbi.nlm.nih.gov/24243024/}}{PMID: 24243024}.
\item[{[21]}] Schaub MA, \textbf{Boyle AP}, Kundaje A, \textsuperscript{\(\dagger\)}Batzoglou S, \textsuperscript{\(\dagger\)}Snyder M. ``Linking disease associations with regulatory information in the human genome.'' \textit{Genome Research} 2012, 22(9):1748--1759. \href{\detokenize{https://pubmed.ncbi.nlm.nih.gov/22955986/}}{PMID: 22955986}.
\item[{[20]}] *Gerstein MB, *Kundaje A, *Hariharan M, *Landt SG, *Yan KK, *Cheng C, *Mu XJ, *Khurana E, *Rozowsky J, *Alexander R, *Min R, *Alves P, Abyzov A, Addleman N, Bhardwaj N, \textbf{Boyle AP}, Cayting P, Charos A, Chen DZ, Cheng Y, Clarke D, Eastman C, Euskirchen G, Frietze S, Fu Y, Gertz J, Grubert F, Harmanci A, Jain P, Kasowski M, Lacroute P, Leng J, Lian J, Monahan H, O'Geen H, Ouyang Z, Partridge EC, Patacsil D, Pauli F, Raha D, Ramirez L, Reddy TE, Reed B, Shi M, Slifer T, Wang J, Wu L, Yang X, Yip KY, Zilberman-Schapira G, Batzoglou S, Sidow A, Farnham PJ, Myers RM, Weissman SM, Snyder M. ``Architecture of the human regulatory network derived from ENCODE data.'' \textit{Nature} 2012, 489(7414):91--100. \href{\detokenize{https://pubmed.ncbi.nlm.nih.gov/22955619/}}{PMID: 22955619}.
\item[{[19]}] The ENCODE Project Consortium. ``An integrated encyclopedia of DNA elements in the human genome.'' \textit{Nature} 2012, 489(7414):57--74. \href{\detokenize{https://pubmed.ncbi.nlm.nih.gov/22955616/}}{PMID: 22955616}.
\item[{[18]}] \textbf{Boyle AP}, Hong EL, Hariharan M, Cheng Y, Schaub MA, Kasowski M, Karczewski KJ, Park J, Hitz BC, Weng S, Cherry JM, Snyder M. ``Annotation of functional variation in personal genomes using RegulomeDB.'' \textit{Genome Research} 2012, 22(9):1790--1797. \href{\detokenize{https://pubmed.ncbi.nlm.nih.gov/22955989/}}{PMID: 22955989}.
\item[{[17]}] *Chen R, *Mias GI, *Li-Pook-Than J, *Jiang L, Lam HYK, Chen R, Miriami E, Karczewski KJ, Hariharan M, Dewey FE, Cheng Y, Clark MJ, Im H, Habegger L, Balasubramanian S, O'Huallachain M, Dudley JT, Hillenmeyer S, Haraksingh R, Sharon D, Euskirchen G, Lacroute P, Bettinger K, \textbf{Boyle AP}, Kasowski M, Grubert F, Seki S, Garcia M, Whirl-Carrillo M, Gallardo M, Blasco MA, Greenberg PL, Snyder P, Klein TE, Altman RB, Butte AJ, Ashley EA, Gerstein M, Nadeau KC, Tang H, Snyder M. ``Personal omics profiling reveals dynamic molecular and medical phenotypes.'' \textit{Cell} 2012, 148(6):1293--1307. \href{\detokenize{https://pubmed.ncbi.nlm.nih.gov/22424236/}}{PMID: 22424236}.
\item[{[16]}] *Song L, *Zhang Z, *Grasfeder LL, *\textbf{Boyle AP}, *Giresi PG, *Lee BK, *Sheffield NC, Graff S, Huss M, Keefe D, Liu Z, London D, McDaniell RM, Shibata Y, Showers KA, Simon JM, Vales T, Wang T, Winter D, Zhang Z, Clarke ND, \textsuperscript{\(\dagger\)}Birney E, \textsuperscript{\(\dagger\)}Iyer VR, \textsuperscript{\(\dagger\)}Crawford GE, \textsuperscript{\(\dagger\)}Lieb JD, \textsuperscript{\(\dagger\)}Furey TS. ``Open chromatin defined by DNaseI and FAIRE identifies regulatory elements that shape cell-type identity.'' \textit{Genome Research} 2011, 21(10):1757--1767. \href{\detokenize{https://pubmed.ncbi.nlm.nih.gov/21750106/}}{PMID: 21750106}.
\item[{[15]}] The ENCODE Project Consortium. ``A User's Guide to the Encyclopedia of DNA Elements (ENCODE).'' \textit{PLoS Biology} 2011, 9(4):e1001046. \href{\detokenize{https://pubmed.ncbi.nlm.nih.gov/21526222/}}{PMID: 21526222}.
\item[{[14]}] \textbf{Boyle AP}, Song L, Lee BK, London D, Keefe D, Birney E, Iyer VR, \textsuperscript{\(\dagger\)}Crawford GE, \textsuperscript{\(\dagger\)}Furey TS. ``High-resolution genome-wide in vivo footprinting of diverse transcription factors in human cells.'' \textit{Genome Research} 2011, 21:456--464. \href{\detokenize{https://pubmed.ncbi.nlm.nih.gov/21106903/}}{PMID: 21106903}.
\item[{[13]}] Georgiev S, \textbf{Boyle AP}, Mukherjee KJAS, Ohler U. ``Evidence-ranked motif identification.'' \textit{Genome Biology} 2010, 11(2):R19. \href{\detokenize{https://pubmed.ncbi.nlm.nih.gov/20156354/}}{PMID: 20156354}.
\item[{[12]}] *Stitzel ML, *Sethupathy P, Pearson DS, Chines PS, Song L, Erdos MR, Welch R, Parker SCJ, \textbf{Boyle AP}, Scott LJ, Margulies EH, Boehnke M, Furey TS, Crawford GE, Collins FS. ``Global epigenomic analysis of primary human pancreatic islets provides insights into type 2 diabetes susceptibility loci.'' \textit{Cell Metabolism} 2010, 12(5):443--455. \href{\detokenize{https://pubmed.ncbi.nlm.nih.gov/21035756/}}{PMID: 21035756}.
\item[{[11]}] McDaniell R, Lee BK, Liu LSAZ, \textbf{Boyle AP}, Erdos MR, Scott LJ, Morken MA, Kucera KS, Battenhouse A, Keefe D, Collins FS, Willard HF, Furey JDLATS, \textsuperscript{\(\dagger\)}Crawford GE, \textsuperscript{\(\dagger\)}Iyer VR, \textsuperscript{\(\dagger\)}Birney E. ``Heritable individual-specific and allele-specific chromatin signatures in humans.'' \textit{Science} 2010, 328(5975):235--239. \href{\detokenize{https://pubmed.ncbi.nlm.nih.gov/20299549/}}{PMID: 20299549}.
\item[{[10]}] Babbitt CC, Fedrigo O, Pfefferle AD, \textbf{Boyle AP}, Horvath JE, Furey TS, Wray GA. ``Both noncoding and protein-coding RNAs contribute to gene expression evolution in the primate brain.'' \textit{Genome Biology and Evolution} 2010, 2:67--79. \href{\detokenize{https://pubmed.ncbi.nlm.nih.gov/20333225/}}{PMID: 20333225}.
\item[{[9]}] Xu X, Tsumagari K, Sowden J, Tawil R, \textbf{Boyle AP}, Song L, Furey TS, Crawford GE, Ehrlich M. ``DNaseI hypersensitivity at gene-poor, FSH dystrophy-linked 4q35.2.'' \textit{Nucleic Acids Research} 2009, 37(22):7381--7393. \href{\detokenize{https://pubmed.ncbi.nlm.nih.gov/19820107/}}{PMID: 19820107}.
\item[{[8]}] \textbf{Boyle AP}, Furey TS. ``High-resolution mapping studies of chromatin and gene regulatory elements.'' \textit{Epigenomics} 2009, 1(2):319--329. \href{\detokenize{https://pubmed.ncbi.nlm.nih.gov/20514362/}}{PMID: 20514362}.
\item[{[7]}] \textbf{Boyle AP}, Guinney J, Crawford GE, Furey TS. ``F-Seq: a feature density estimator for high-throughput sequence tags.'' \textit{Bioinformatics} 2008, 24(21):2537--2538. \href{\detokenize{https://pubmed.ncbi.nlm.nih.gov/18784119/}}{PMID: 18784119}.
\item[{[6]}] \textbf{Boyle AP}, Davis S, Shulha HP, Meltzer P, Margulies EH, Weng Z, \textsuperscript{\(\dagger\)}Furey TS, \textsuperscript{\(\dagger\)}Crawford GE. ``High-resolution mapping and characterization of open chromatin across the genome.'' \textit{Cell} 2008, 132(2):311--322. \href{\detokenize{https://pubmed.ncbi.nlm.nih.gov/18243105/}}{PMID: 18243105}.
\item[{[5]}] \textbf{Boyle AP}, Boyle JA, Bridges SM. ``Identification of Regulatory Elements in Archaea using Self-Organizing Maps.'' \textit{Proc RECOMB} 2004. \href{\detokenize{http://boylelab.org/pubs/RECOMB_2004_Boyle.pdf}}{Publication link}.
\item[{[4]}] \textbf{Boyle AP}, Boyle JA. ``Global analysis of microbial translation initiation regions.'' \textit{Journal of the Mississippi Academy of Sciences} 2003, 48(3):138--150. \href{\detokenize{http://www.thefreelibrary.com/Global+analysis+of+microbial+translation+initiation+regions-a0105160258}}{Publication link}.
\item[{[3]}] \textbf{Boyle AP}, Bridges S. ``Clustering of archael gene regulatory regions.'' \textit{FASEB Journal} 2003, 17(5):A985-A985.
\item[{[2]}] \textbf{Boyle AP}, Boyle JA. ``Visualization of aligned genomic open reading frame data.'' \textit{Biochemistry and Molecular Biology Education} 2003, 31(1):64--68. \href{\detokenize{https://doi.org/10.1002/bmb.2003.494031010144}}{doi: 10.1002/bmb.2003.494031010144}.
\item[{[1]}] Wan X, Boyle JA, Bridges SM, \textbf{Boyle AP}. ``Interactive clustering for exploration of genomic data.'' \textit{Proceedings of the Artificial Neural Networks in Engineering Conference} 2002, 12:753--758. \href{\detokenize{http://boylelab.org/pubs/ANNIE_2002_Wan.pdf}}{Publication link}.
\end{enumerate}


% Generated by scripts/build_cv.py. Do not edit directly.
% Source: cv/patents.bib.
\cvsection{Patents}
\begin{enumerate}[leftmargin=2.7em,labelsep=0.55em,itemsep=0.65em,parsep=0pt,topsep=0.25em]
\item[{[1]}] Karczewski K, Snyder M, Butte AJ, Dudley JT, Hong E, \textbf{Boyle A}, Cherry MJ, Park J. ``Method and System for the Use of Biomarkers for Regulatory Dysfunction in Disease.'' Granted United States patent 9,946,835, 2018. \href{\detokenize{https://patents.google.com/patent/US9946835B2/en}}{Patent record}.
\end{enumerate}


\end{document}

